\section{Problem Summary}
Given an integer array, return all distinct triplets whose sum is zero.

The difficulty is not only finding triplets fast enough. You also need to avoid duplicate answers when the input contains repeated values.

\section{Recognition Pattern}
This is a \textbf{sorting + two pointers} problem.

In interviews, recognize it when:
\begin{itemize}
\item you need unique pairs or triplets
\item the order of the array does not matter
\item sorting lets you reason about sums moving left or right
\item duplicate handling is part of the difficulty
\end{itemize}

\section{Key Idea}
Sort the array first. Then fix one index \(i\) and solve a two-sum problem on the suffix \([i + 1, n - 1]\) with two pointers.

Because the array is sorted:
\begin{itemize}
\item if the current sum is too small, move the left pointer right
\item if the current sum is too large, move the right pointer left
\item after finding a valid triplet, skip duplicates on both sides
\end{itemize}

Sorting turns duplicate control from a bookkeeping mess into a local pointer rule.

\section{Step-by-step Reasoning}
The direct brute-force solution tries every triple, which costs \(O(n^3)\).

A better but still clumsy idea is to fix one number and use a hash set for the remaining two-sum search. That reduces the time to \(O(n^2)\), but keeping the output unique becomes awkward.

Sorting improves both issues at once. Once the numbers are ordered, moving a pointer has a predictable effect on the sum, and duplicate values sit next to each other, so they are easy to skip.

That is the key transition: instead of treating uniqueness as a separate cleanup pass, build it into the scan itself.

\section{Complexity}
\begin{itemize}
\item Time: \(O(n^2)\) overall
\item Space: \(O(1)\) auxiliary, ignoring the output and the sort implementation details
\end{itemize}

\section{Common Pitfalls}
\begin{itemize}
\item Forgetting to skip duplicate anchors creates repeated triplets.
\item After finding one triplet, failing to skip duplicate left or right values also creates repeats.
\item Using the two-pointer pattern on an unsorted array does not work.
\item Missing the early stop when the anchor becomes positive wastes time, because three positive numbers cannot sum to zero in a sorted array.
\end{itemize}

\section{Variants / Follow-ups}
This pattern generalizes directly to \textbf{3Sum Closest}, \textbf{4Sum}, and the broader \(k\)-sum family.

The main reusable lesson is that sorting is often worth it when uniqueness and pointer movement need to be handled at the same time.
